\section{Other index considerations}
\label{sec:index-operational-considerations}

Indexes can automatically update themselves on index, however currently enabled by default is grouped MPF index, which due to bugs cannot correctly 
return minimizers which leads to recomputing them when no explicit index read policy is specified and it doesn't support update on insert.
More details are presented in the Appendix 1, where I also present concrete examples of how the BioseqDB should be used.

This bug was caused by the fact that grouped MPF not only produces relational index, but also saves that compressed relation chunks into separate file that is used to enable faster reads during JOINs.
The reader by default rejects ordinary single table queries, so you need to change the GUC to disable grouped-MPF (and use regular MPF instead).
Saving file enables the JOIN to concurrently access stored data which is a cheap workaround around the decision to only use single Postgres reader process.

\begin{figure}[H]
\centering
\begin{tikzpicture}[x=1cm,y=1cm,font=\footnotesize]
  \foreach \shift/\mode/\heading in
    {0/0/{(a) Ordinary MPF},7/1/{(b) Grouped MPF}} {
    \begin{scope}[xshift=\shift cm]
      \node[anchor=west,font=\bfseries] at (.1,10.65) {\heading};
      \node at (4.1,10.15) {\texttt{refs}: targets};
      \draw[fill=black!4,draw=black!55] (2.4,8.08) rectangle (5.8,9.88);
      \draw[fill=indexGreen!16,draw=indexGreen] (2.4,8.68) rectangle (5.8,9.28);
      \draw[black!55] (2.4,9.28)--(5.8,9.28);
      \draw[black!55] (3.45,8.08)--(3.45,9.88);
      \node at (2.925,9.58) {id};
      \node at (4.625,9.58) {seq};
      \node at (2.925,8.98) {$t_0$};
      \node at (4.625,8.98) {$\cdots$};
      \node[text=black!50] at (2.925,8.38) {$t_1$};
      \node[text=black!50] at (4.625,8.38) {$\cdots$};
      \draw[index arrow] (4.1,7.98)--(4.1,7.27);
      \node[anchor=west] at (4.3,7.68) {index for $t_0$};

      \ifnum\mode=0
        \node at (4.1,7.04) {MPF index relation};
        \foreach \yy in {6.6,6.23,5.86} {
          \draw[fill=black!6,draw=black!55] (2.7,\yy-.12) rectangle (5.5,\yy+.12);
          \draw[fill=indexOrange!25,draw=indexOrange] (3.25,\yy-.12) rectangle (5.5,\yy+.12);
        }
        \draw[index arrow,indexOrange] (4.1,5.59)--(4.1,4.75);
        \node[anchor=west,align=left] at (4.32,5.18) {load and\\assemble};
        \draw[fill=indexOrange!10,draw=indexOrange] (3.0,2.25) rectangle (6.15,4.65);
        \node at (4.575,4.3) {private CPU lookup};
      \else
        \node at (4.1,7.04) {MPF reader image};
        \draw[fill=indexGreen!7,draw=indexGreen]
          (2.7,5.74)--(5.5,5.74)--(5.5,6.43)--(5.21,6.72)--(2.7,6.72)--cycle;
        \draw[indexGreen] (5.21,6.72)--(5.21,6.43)--(5.5,6.43);
        \foreach \yy in {6.47,6.2,5.93}
          \draw[line width=3pt,indexGreen!45] (2.98,\yy)--(4.92,\yy);
        \draw[index arrow,indexGreen] (4.1,5.59)--(4.1,4.75);
        \node[anchor=west,align=left] at (4.32,5.18) {open and\\\texttt{mmap}};
        \draw[dashed,fill=indexGreen!6,draw=indexGreen] (3.0,2.25) rectangle (6.15,4.65);
        \node at (4.575,4.3) {mapped array view};
      \fi
      \node[align=center] at (4.575,3.36) {one reader\\for target $t_0$};

      \node at (1.0,5.02) {\texttt{queries}};
      \draw[fill=indexBlue!6,draw=indexBlue] (.1,2.25) rectangle (1.9,4.65);
      \draw[indexBlue] (.75,2.25)--(.75,4.65);
      \foreach \yy in {4.05,3.45,2.85}
        \draw[indexBlue] (.1,\yy)--(1.9,\yy);
      \node at (.425,4.35) {id};
      \node at (1.325,4.35) {seq};
      \foreach \yy/\q in {3.75/0,3.15/1,2.55/2} {
        \node at (.425,\yy) {$q_{\q}$};
        \node at (1.325,\yy) {$\cdots$};
        \draw[index arrow,indexBlue] (1.98,\yy)--(2.9,\yy);
      }
      \node[align=center] at (1.0,1.72) {keep query\\rows};
      \draw[index arrow] (4.575,2.14)--(4.575,1.55);
      \node at (4.575,1.27) {target positions for anchors};
      \draw[black!45] (3.1,.8)--(6.05,.8);
      \foreach \xx in {3.32,4.02,5.7}
        \fill[indexGreen] (\xx,.8) circle (2.2pt);
    \end{scope}
  }
\end{tikzpicture}
\caption{For large joins grouped-MPF uses saved direct binary dump of relational index, which speeds up the alignment by reducing time to sequentially read the backing table}
\label{fig:index-join-read-paths}
\end{figure}

All indexes in BioseqDB can be automatically updated on insert (but as noted earlier, grouped-MPF will produce error on insert due to buggy implementation)
\begin{table}[htbp]
\centering
\small
\caption{Maintenance policies}
\label{tab:index-update-policy}
\begin{tabularx}{\linewidth}{@{}p{.31\linewidth}X@{}}
\toprule
\textbf{Policy} & \textbf{Description}\\
\midrule
\texttt{UPDATE\_ON\_INSERT} & When user inserts data, we automatically compute and store minimizers for new sequence \\
\texttt{UPDATE\_MANUALLY} & User needs to manually call \texttt{bioseq.update\_index} \\
\bottomrule
\end{tabularx}
\end{table}

\begin{table}[htbp]
\centering
\small
\caption{Default policies in BioseqDB}
\label{tab:index-policy-defaults}
\begin{tabularx}{\linewidth}{@{}p{.34\linewidth}p{.30\linewidth}X@{}}
\toprule
\textbf{Index} & \textbf{Policy for user updates} & \textbf{Description}\\
\midrule
\path{minimizers_per_fragment} & \texttt{UPDATE\_ON\_INSERT} & Default minimizer index\\
\path{minimizers_per_sequence} & \texttt{UPDATE\_MANUALLY} & Not selected by default\\
\path{symbol_counts} & \texttt{UPDATE\_ON\_INSERT} & Used when user demans to count symbols in a sequence\\
\bottomrule
\end{tabularx}
\end{table}

